\chapter{维蒂希反应及其他形成 C=C 的反应}\label{ch:b2:wittig}

雌性家蝇以一种含 23 个碳、带有一个双键的烃吸引雄蝇，该双键位于
碳 9 与碳 10 之间，构型为 (\textit{Z})。要大量制备它，化学家必须把这个双键
恰好放在该处并保持该构型：消除反应会使双键分散到几个位置，且同时得到两种构型。
本章换一种方式构造双键，即把两个片段连接在一个羰基上，
使双键恰好出现在原来 \ce{C=O} 所在的位置。
% ledger: use:muscalure

\begin{recall}
\Cref{ch:b2:enolates-aldol}：羟醛缩合。\Cref{ch:b2:alkene-redox}：炔烃的
催化氢化为 (\textit{Z})-烯烃，所用为中毒催化剂。本科一年级分册：E1、E2 与
札依采夫规则、双分子取代、格氏试剂与极性反转，以及
(\textit{Z})/(\textit{E}) 构型标记。第 1 册：原子经济性。
\end{recall}

\section{鏻盐与磷叶立德}

\begin{definition}[磷叶立德]\label{def:b2:wittig:ylide}
\emph{鏻盐}\index{鏻盐} $\mathrm{R_3P^+{-}CH_2R'\ X^-}$ 可由膦与卤代烷经
双分子取代制得，常用膦为三苯基膦 \ce{Ph3P}。去掉与磷相连的碳上的
一个氢，便得到\emph{磷叶立德}\index{磷叶立德}，
即一种中性分子，两相邻原子带相反电荷，$\mathrm{Ph_3P^+{-}CH^-{-}R'}$，也可写作 $\mathrm{Ph_3P{=}CHR'}$。
\emph{稳定叶立德}\index{稳定叶立德}在该碳上连有吸电子基团，
（酯基、酮基、腈基），使负电荷发生离域。
\end{definition}

\begin{omfigure}
\setchemfig{atom sep=1.9em}
\schemestart
\chemfig{Ph_3P} \+ \chemfig{CH_3-I}
\arrow{->[S$_\mathrm{N}$2]}[,0.9]
\chemfig{Ph_3P^{+}-CH_3}\;\chemfig{I^{-}}
\arrow{->[\ce{BuLi}][$-$ \ce{BuH}, \ce{LiI}]}[,1.3]
\chemfig{Ph_3P^{+}-\chemabove{C}{\scriptstyle -}H_2}
\arrow{<->}[,0.8]
\chemfig{Ph_3P=CH_2}
\schemestop
\par\vspace{6pt}

{\small 亚甲基三苯基磷烷，即最简单的叶立德。三苯基膦从碘甲烷中
取代出碘离子；丁基锂夺去甲基上的一个质子。右侧两个结构是
同一分子的两种写法，其中电荷分离的写法更接近实际情况。}
\end{omfigure}

\begin{proposition}[制备叶立德]\label{prop:b2:wittig:ylide-acidity}
鏻盐在与磷相邻的碳上的酸性远强于烷烃。简单的烷基
鏻盐仍需很强的碱（丁基锂、氢化钠、氨基钠），并须在干燥、
惰性条件下进行；若盐中 \ce{CH2} 上还连有羰基，则弱碱（碳酸钠、
氢氧化钠）即可使其去质子，所生成的稳定叶立德可以分离并保存。
\end{proposition}

\begin{proof}[论证]
与碳负离子相邻的带正电磷原子，通过其电荷以及大体积磷原子的
极化率使碳负离子稳定：简单叶立德是一种稳定的碳负离子，但稳定程度
不及烯醇负离子，因此需要强碱。相邻的羰基可将负电荷分散到氧上，与烯醇负离子相同，
（见\cref{prop:b2:enolates-aldol:alpha-acidity}），在磷的作用之外再叠加，使共轭酸
的酸性强到弱碱即可使之去质子；而叶立德碱性弱得多，不再与水或空气反应。
\end{proof}

\section{维蒂希反应}

\begin{definition}[维蒂希反应]\label{def:b2:wittig:wittig}
磷叶立德与醛或酮发生的\emph{维蒂希反应}\index{维蒂希反应}，
生成烯烃和三苯基氧化膦：
\begin{equation*}
\mathrm{Ph_3P{=}CHR + R'CHO \longrightarrow R'CH{=}CHR + Ph_3P{=}O}.
\end{equation*}
反应经由\emph{氧磷杂环丁烷}\index{氧磷杂环丁烷}进行，它是由磷、两个
碳和氧组成的四元环，由叶立德碳与羰基碳成键，同时羰基氧
与磷成键而形成。
\end{definition}

\begin{omfigure}
\setchemfig{atom sep=1.9em}
\schemestart
\chemfig{Ph_3P=CH-R} \+ \chemfig{O=CH-R'}
\arrow{->[加成]}[,1.0]
\chemfig{Ph_3P?-[6]CH(-[4]R)-[0]CH(-[0]R')-[2]O?}
\schemestop

\medskip
\schemestart
\arrow{->[消除]}[,1.1]
\chemfig{R-CH=CH-R'} \+ \chemfig{Ph_3P=O}
\schemestop
\par\vspace{6pt}

{\small 维蒂希反应示意。叶立德碳与羰基碳成键，同时氧与磷
成键，一步完成，生成氧磷杂环丁烷。随后该环朝另一方向开环，断开
\ce{P-C} 与 \ce{C-O} 键：\ce{C=C} 键形成于原先参与反应的两个碳之间，同时
三苯基氧化膦的 \ce{P=O} 键也随之形成。}
\end{omfigure}

\begin{proposition}[双键恰好在位]\label{prop:b2:wittig:regiospecific}
在维蒂希反应中，新生成的 \ce{C=C} 键连接原来的羰基碳与原来的叶立德碳，
且不会出现在其他位置。
\end{proposition}

\begin{proof}[论证]
仅有的成键与断键都发生在四元环上：叶立德碳与羰基碳成键，
随后断开 \ce{C-P} 与 \ce{C-O} 键。没有氢原子迁移，沿碳链也不生成碳正离子或碳负离子，
因此双键不会迁移。相反，消除反应需要在相邻若干个碳上的氢之间进行选择
（札依采夫规则），且 E1 反应可能发生重排。
\end{proof}

\begin{proposition}[推动力]\label{prop:b2:wittig:driving}
三苯基氧化膦中 \ce{P=O} 键的形成，使维蒂希反应
极为有利，且是不可逆的。
\end{proposition}

\begin{proof}[论证]
反应中，一个 \ce{C=O} 键与一个 \ce{P-C} 键（叶立德中的 \ce{P=C} 主要为
\ce{P+-C-} 单键）被替换为一个 \ce{C=C} 键与一个 \ce{P=O} 键。氧化膦的 \ce{P=O} 键
非常牢固，远强于所断裂的 \ce{P-C} 键；所失去的 \ce{C=O} 键部分地由新形成的
\ce{C=C} 键补偿，平衡显然偏向生成物一侧。磷是本
反应的氧“汇”。
\end{proof}

\section{立体选择性}

\begin{proposition}[烯烃的几何构型]\label{prop:b2:wittig:stereo}
与醛反应时，非稳定叶立德（\ce{Ph3P=CHR}，R 为烷基）主要生成 (\textit{Z})-烯烃，
在无盐条件下。稳定叶立德（\ce{Ph3P=CH-COOR}，\ce{Ph3P=CH-COR}）主要生成
(\textit{E})-烯烃。仅由芳基稳定的叶立德则得到混合物。
\end{proposition}

\admitted

以上均为观察结果。通常的解释把维蒂希反应看作
速率与稳定性之间的竞争。对于活泼的非稳定叶立德，氧磷杂环丁烷迅速且不可逆地生成，
经由位阻较小的过渡态，两个取代基最终位于
环的同一侧；随后在保持这一构型的情况下分解，受动力学控制得到
(\textit{Z})-烯烃。对于稳定叶立德，第一步较慢且可逆；
氧磷杂环丁烷通过原料相互转化，其中取代基位于
两侧、张力较小的那一种占优，得到 (\textit{E})-烯烃，受热力学控制。

\begin{omfigure}
\setchemfig{atom sep=1.9em}
\schemestart
\chemfig{Ph_3P=CH-CH_2CH_3} \+ \chemfig{O=CH-Ph}
\arrow{->}[,0.8]
\chemfig{Ph-[:-30]=[:30]-[:90]CH_2CH_3}
\schemestop

\medskip
\schemestart
\chemfig{Ph_3P=CH-COOEt} \+ \chemfig{O=CH-Ph}
\arrow{->}[,0.8]
\chemfig{Ph-[:-30]=[:30]-[:-30]COOEt}
\schemestop
\par\vspace{6pt}

{\small 两种叶立德，一种醛。上：非稳定的丙叉叶立德主要生成
(\textit{Z})-1-苯基丁-1-烯。下：稳定叶立德主要生成乙基
(\textit{E})-3-苯基丙烯酸酯（肉桂酸乙酯）。两者中均生成的三苯基氧化膦未
画出。}
\end{omfigure}

\section{霍纳–沃兹沃思–埃蒙斯反应}

\begin{definition}[膦酸酯与 HWE 反应]\label{def:b2:wittig:hwe}
称\emph{膦酸酯}\index{膦酸酯}为膦酸的酯，其一般式为 $\mathrm{(RO)_2P({=}O){-}R'}$，
含 \ce{P-C} 键。它由亚磷酸三烷基酯与卤代烷经米歇利斯–阿尔布佐夫反应制得，
$\mathrm{P(OEt)_3 + R'Br \to (EtO)_2P(O)R' + EtBr}$。在
\emph{霍纳–沃兹沃思–埃蒙斯反应}\index{霍纳–沃兹沃思–埃蒙斯反应}（HWE）中，膦酸酯的负离子
（其 \ce{CH2} 上带吸电子基团）与醛加成，并生成
烯烃与二烷基磷酸根负离子。
\end{definition}

\begin{omfigure}
\setchemfig{atom sep=1.8em}
\schemestart
\chemfig{EtO-P(=[2]O)(-[6]OEt)-CH_2-COOEt}
\arrow{->[\ce{NaH}][$-$ \ce{H2}]}[,0.9]
\chemfig{EtO-P(=[2]O)(-[6]OEt)-\chemabove{C}{\scriptstyle -}H-COOEt}
\schemestop

\medskip
\schemestart
\arrow{->[\ce{PhCHO}]}[,0.9]
\chemfig{*6(=-=(-[:30]=[:-30]-[:30](=[:90]O)-[:-30]O-[:30]Et)-=-)}
\+
\chemfig{EtO-P(=[2]O)(-[6]O^{-})-OEt}
\schemestop
\par\vspace{6pt}

{\small 膦酰乙酸三乙酯与苯甲醛的霍纳–沃兹沃思–埃蒙斯反应。氢化钠
夺去膦酸酯与酯基之间的氢；负离子加成至醛上，
四元环中间体如维蒂希反应那样分解。产物为
(\textit{E})-肉桂酸乙酯；副产物为二乙基磷酸根负离子，易溶于水。}
\end{omfigure}

\begin{proposition}[化学家青睐 HWE 反应的原因]\label{prop:b2:wittig:hwe}
HWE 反应主要生成 (\textit{E})-烯烃，其含磷副产物为水溶性
盐，可用简单的水洗除去；膦酸酯负离子的亲核性也比
相应的稳定叶立德更强，它还能与酮反应。
\end{proposition}

\begin{proof}[论证]
膦酸酯负离子是稳定的碳负离子，因此其加成是可逆的，反应
受热力学控制，得到 (\textit{E})-产物，与稳定叶立德相同。与磷叶立德不同，它
带有完整负电荷，无相邻阳离子加以补偿，因此反应活性更高。其
副产物 \ce{(EtO)2PO2-} 为离子型盐，而三苯基氧化膦是中性有机
固体，需通过色谱法或结晶与产物分离。
\end{proof}

\section{方法的选择}

\begin{method}[维蒂希切断]\label{met:b2:wittig:wittig-disconnection}
\begin{enumerate}
  \item 切断目标分子中的 \ce{C=C} 键；一个碳成为羰基碳，另一个成为叶立德碳。
        （其前体为鏻盐，由卤化物制得）。
  \item 两种做法中，优先选择卤化物为伯卤代烷的那一种（便于 
        \ce{Ph3P} 双分子取代），且羰基化合物为醛的那一种。
  \item 按所需几何构型选择试剂：得到 (\textit{Z})-构型用非稳定叶立德，
        得到 (\textit{E})-构型则用稳定叶立德或膦酸酯（HWE）。
  \item 检查分子中其他基团是否会与所用的碱反应。
\end{enumerate}
\end{method}

\begin{method}[选择生成 C=C 的反应]\label{met:b2:wittig:choose-cc}
\begin{center}
\small
\begin{tabular}{@{}lll@{}}
\hline
方法 & \ce{C=C} 的位置 & 构型 \\
\hline
E2 或 E1 消除 & 札依采夫规则，混合物 & 主要为 (\textit{E})，混合物 \\
羟醛缩合 & 与 \ce{C=O} 共轭 & (\textit{E}) \\
维蒂希反应，非稳定叶立德 & 原 \ce{C=O} 处 & (\textit{Z}) \\
维蒂希反应，稳定叶立德；HWE & 原 \ce{C=O} 处 & (\textit{E}) \\
炔烃，再用中毒催化剂 & 原三键处 & (\textit{Z}) \\
\hline
\end{tabular}
\end{center}
\medskip
在 \ce{C=C} 处连接两个大小相近的片段时，维蒂希反应及其同类反应通常
是首选；烯烃复分解可交换两个烯烃的端基，将在本科
三年级分册中介绍。
\end{method}

\begin{history}[成为试剂的叶立德]
维蒂希在 1950 年代初研究磷叶立德时发现，亚甲基三苯基磷烷
可将二苯甲酮转化为 1,1-二苯乙烯和三苯基氧化膦；他与学生乌尔里希·舍尔科普夫
一道把这一观察发展为通用方法，时间是 1954 年。工业界迅速采用了该方法，
用于维生素 A 等物质的生产，维蒂希于 1979 年共享诺贝尔化学奖。
\end{history}

\begin{safety}{\ghs[5mm]{GHS02}\,\ghs[5mm]{GHS05}\,\ghs[5mm]{GHS07}\,\ghs[5mm]{GHS08}}
丁基锂溶液易燃、有腐蚀性，且遇水剧烈反应；氢化钠
遇水会放出氢气。三苯基膦有害且对眼睛有腐蚀性，亚磷酸三乙酯
易燃且有害；三苯基氧化膦有害。
% ledger: ghs:BuLi, ghs:PPh3, ghs:triethyl-phosphite, ghs:Ph3PO
\end{safety}

\section{练习}

\begin{exercise}[$\star$]\label{exo:b2:wittig:1}
写出由三苯基膦与卤代烷制备叶立德 \ce{Ph3P=CHCH3} 的过程。
\end{exercise}

\begin{exercise}[$\star$]\label{exo:b2:wittig:2}
求苯甲醛与由碘甲烷制得的叶立德进行维蒂希反应的产物。
\end{exercise}

\begin{exercise}[$\star$]\label{exo:b2:wittig:3}
将下列叶立德分类为稳定叶立德、非稳定叶立德或芳基稳定叶立德：\ce{Ph3P=CHCOOEt}、\ce{Ph3P=CHCH3}、
\ce{Ph3P=CHCOCH3}、\ce{Ph3P=CHC6H5}。哪一种需要用丁基锂？
\end{exercise}

\begin{exercise}[$\star$]\label{exo:b2:wittig:4}
求膦酰乙酸三乙酯、氢化钠与丙醛反应的产物，并给出其构型。
\end{exercise}

\begin{exercise}[$\star\star$]\label{exo:b2:wittig:5}
需制备亚甲基环己烷。比较由环己酮出发的维蒂希路线，与
1-甲基环己-1-醇的酸催化脱水路线。
\end{exercise}

\begin{exercise}[$\star\star$]\label{exo:b2:wittig:6}
对 (\textit{E})-1,2-二苯乙烯（芪）作维蒂希切断。为什么维蒂希
路线不适合制备纯的 (\textit{E})-芪？应改用什么方法？
\end{exercise}

\begin{exercise}[$\star\star$]\label{exo:b2:wittig:7}
丙醛分别与 \ce{Ph3P=CHCH3} 和 \ce{Ph3P=CHCOOEt} 反应。求
出各产物及其主要构型。
\end{exercise}

\begin{exercise}[$\star\star$]\label{exo:b2:wittig:8}
计算用 \ce{Ph3P=CH2} 使环己酮亚甲基化时的原子经济性。
\end{exercise}

\begin{exercise}[$\star\star$]\label{exo:b2:wittig:9}
写出亚磷酸三乙酯与溴乙酸乙酯发生的米歇利斯–阿尔布佐夫反应，
分两步进行：先是磷的取代，再是溴离子使其脱烷基。
\end{exercise}

\begin{exercise}[$\star\star\star$]\label{exo:b2:wittig:10}
提出以 (\textit{E})-丁-2-烯二醛
\ce{OHC-CH=CH-CHO} 为原料，经两次维蒂希反应合成 1,6-二苯基己-1,3,5-三烯的路线，并说明所用叶立德及其用量。
\end{exercise}

\begin{exercise}[$\star\star\star$]\label{exo:b2:wittig:11}
(\textit{E})-4-苯基丁-3-烯-2-酮可由羟醛缩合或维蒂希反应制得。
写出两条路线并加以比较。
\end{exercise}

\begin{exercise}[$\star\star\star$]\label{exo:b2:wittig:12}
提出两条合成 (\textit{Z})-辛-4-烯的路线，一条用维蒂希反应，另一条经由炔烃，
并说明各路线分别留下哪些副产物。
\end{exercise}

\section{问题：用维蒂希反应合成信息素}

\begin{problem}[{周末问题——(Z)-烯烃的逆合成分析，非稳定叶立德及其选择性，
三苯基氧化膦的代价，以及炔烃替代方案}]
\label{pb:b2:wittig:1}
目标物为 (\textit{Z})-二十三碳-9-烯，\ce{CH3(CH2)7CH=CH(CH2)12CH3}，即
开头所述的家蝇信息素。摩尔质量（\unit{g/mol}）：C 12.011，H 1.008，O 15.999，P 30.974，Br 79.904。
% ledger: use:muscalure, aw:C, aw:H, aw:O, aw:P, aw:Br, ghs:nonanal, ghs:BuLi

\medskip
\noindent\textbf{第一部分——逆合成分析。}
\begin{enumerate}
  \item 写出目标分子的分子式，并判断它是否为烯烃。
  \item 维蒂希切断断开的是哪一个键？
  \item 给出两组可能的配对（羰基化合物、鏻盐）。
  \item 对于壬醛与叶立德这一组，由哪种卤代烷制得该鏻盐？
  \item 为什么伯卤代烷适合用于这一步？
  \item 写出鏻盐的生成反应。
\end{enumerate}

\medskip
\noindent\textbf{第二部分——叶立德与几何构型。}
\begin{enumerate}[resume]
  \item 哪种碱使该盐去质子化？为什么氢氧化钠不行？
  \item 该叶立德是否为稳定叶立德？这预示着怎样的几何构型？
  \item 写出壬醛与之反应生成的氧磷杂环丁烷。
  \item 为什么反应必须在无水及惰性气体中进行？
  \item 为什么双键恰好位于第 9 和第 10 号碳之间？
  \item 十三碳-9-醇经酸催化脱水会得到什么？
\end{enumerate}

\medskip
\noindent\textbf{第三部分——三苯基膦氧化物。}
\begin{enumerate}[resume]
  \item 写出叶立德与壬醛反应的总方程式。
  \item 计算目标分子、壬醛、鏻盐和
        三苯基膦氧化物的摩尔质量。
  \item 计算维蒂希步骤（叶立德加醛）的原子经济性。
  \item 如何将三苯基膦氧化物与烃类分离？
  \item 为什么生成三苯基膦氧化物是驱动力？
  \item 为什么霍纳–沃兹沃思–埃蒙斯反应不适合该目标分子？
\end{enumerate}

\medskip
\noindent\textbf{第四部分——炔烃路线。}
\begin{enumerate}[resume]
  \item 哪种炔烃经加氢可得到目标分子？使用何种催化剂？
  \item 由癸-1-炔与卤代烷构筑该炔烃：用哪种卤化物、哪种碱？
  \item 为什么加氢必须停留在烯烃阶段？催化剂如何保证这一点？
  \item 比较两条路线：步骤数、几何构型控制、副产物。
  \item 加氢步骤的原子经济性是多少？
  \item 说明经维蒂希路线每克信息素
        同时生成的三苯基膦氧化物的质量。
\end{enumerate}
\end{problem}
